% §7 Exercises
%
% Cross-reference resolution status:
%
% DEFINED here:
%   \label{ex:ch7-1} through \label{ex:ch7-11}
%
% RESOLVED (used here, defined elsewhere):
%   \Cref{prop:binomial-unbiased}         -- §7.2 (Plan 7 Task 3)
%   \Cref{prop:binomial-variance}         -- §7.2 (Plan 7 Task 3)
%   \Cref{prop:binomial-mle}              -- §7.2 (Plan 7 Task 3)
%   \Cref{def:wald-ci}                    -- §7.3 (Plan 7 Task 4)
%   \Cref{def:wilson-ci}                  -- §7.3 (Plan 7 Task 4)
%   \Cref{def:clopper-pearson-ci}         -- §7.3 (Plan 7 Task 4)
%   \Cref{def:jeffreys-ci}                -- §7.3 (Plan 7 Task 4)
%   \Cref{def:ppv}                        -- §7.4 (Plan 7 Task 5)
%   \Cref{def:npv}                        -- §7.4 (Plan 7 Task 5)
%   \Cref{def:f1}                         -- §7.4 (Plan 7 Task 5)
%   \Cref{thm:delta-method}               -- §7.4 (Plan 7 Task 5)
%   \Cref{prop:ppv-variance}              -- §7.4 (Plan 7 Task 5)
%
% Difficulty-marker macros (not in alex.sty; defined locally if not yet defined):
\providecommand{\exstar}{$\star$}
\providecommand{\exstarstar}{$\star\star$}
\providecommand{\exstarstarstar}{$\star\star\star$}

\section*{Exercises}
\addcontentsline{toc}{section}{Exercises}

\begin{enumerate}

% -----------------------------------------------------------------------
% WARM-UP  (★ × 3): exercises 1, 2, 3
% -----------------------------------------------------------------------

\item \exstar
\label{ex:ch7-1}
\textbf{Compute PPV.}
A classifier has rates $\fprate = 0.05$ and $\fnrate = 0.10$, and is
deployed on three workloads whose prevalences are $\pi \in \{0.01,
0.1, 0.5\}$.
\begin{enumerate}
  \item[(a)] Substitute each prevalence into \Cref{def:ppv} and
    compute the positive predictive value. Report to four decimal
    places.
  \item[(b)] Comment on the magnitude of $\PPV$ at the three
    prevalences. The classifier's $(\fprate, \fnrate)$ has not
    moved across the three workloads; only $\pi$ has.
  \item[(c)] In one sentence, name which of the three workloads
    makes the classifier appear most precise and which makes it
    appear least precise, and identify the §7.4 paragraph that
    flags this prevalence-dependence as the central caveat for
    reporting precision in isolation.
\end{enumerate}

\item \exstar
\label{ex:ch7-2}
\textbf{Compute $\Fone$.}
The same classifier $\fprate = 0.05, \fnrate = 0.10$ is deployed on
the balanced workload $\pi = 0.5$.
\begin{enumerate}
  \item[(a)] Compute $\Recall = 1 - \fnrate$.
  \item[(b)] Compute $\Precision = \PPV$ at $\pi = 0.5$ by
    \Cref{def:ppv}.
  \item[(c)] Substitute the two values into \Cref{def:f1} to compute
    $\Fone$. Report to four decimal places.
\end{enumerate}

\item \exstar
\label{ex:ch7-3}
\textbf{Wilson interval evaluation.}
A classifier was run on $n_0 = 100$ truth-$\bot$ trials and produced
$\mathrm{FP} = 10$ false positives.
\begin{enumerate}
  \item[(a)] Compute the empirical rate
    $\hat{\fprate} = \mathrm{FP}/n_0$.
  \item[(b)] Apply \Cref{def:wilson-ci} at $z_{0.025} = 1.96$ to
    compute the Wilson 95\% confidence interval for $\fprate$.
    Report the center, the half-width, and the resulting interval
    endpoints to four decimal places.
  \item[(c)] Compare to the Wald interval \Cref{def:wald-ci} at the
    same data: by how much does the Wilson center differ from
    $\hat{\fprate}$, and in which direction is it shifted? Refer
    back to §7.3 for the score-statistic shift toward $1/2$.
\end{enumerate}

% -----------------------------------------------------------------------
% MEDIUM  (★★ × 5): exercises 4, 5, 6, 7, 8
% -----------------------------------------------------------------------

\item \exstarstar
\label{ex:ch7-4}
\textbf{Notebook-tied CI comparison.}
Open the chapter 7 companion notebook
\texttt{python/notebooks/ch07-classification-measures.ipynb}.
\begin{enumerate}
  \item[(a)] For each combination $n_0 \in \{20, 100, 500\}$ and
    $\fprate \in \{0.01, 0.05, 0.5\}$, simulate $M = 10000$
    independent draws of $\mathrm{FP} \sim
    \mathrm{Binomial}(n_0, \fprate)$, construct the four intervals
    of \Cref{def:wald-ci,def:wilson-ci,def:clopper-pearson-ci,def:jeffreys-ci}
    at each draw, and tabulate the empirical coverage (fraction
    of intervals containing the true $\fprate$).
  \item[(b)] From the resulting $3 \times 3$ table for each
    interval, identify the regime $(n_0, \fprate)$ where Wald
    coverage falls farthest below the nominal $0.95$.
  \item[(c)] In one sentence, state which of the three alternative
    intervals recovers most of the lost coverage in that regime
    and which has the highest coverage overall. Refer to
    \Cref{def:wilson-ci,def:clopper-pearson-ci,def:jeffreys-ci}
    by name.
\end{enumerate}

\item \exstarstar
\label{ex:ch7-5}
\textbf{Delta method for log-odds.}
The log-odds transformation $g(\fprate) = \log[\fprate / (1 -
\fprate)]$ maps the binomial parameter from $(0, 1)$ to the entire
real line and is the standard re-parametrization in logistic
regression.
\begin{enumerate}
  \item[(a)] Differentiate $g$ and write $g'(\fprate)$ as a function
    of $\fprate$. Simplify the two-term derivative to a single
    rational expression.
  \item[(b)] Apply \Cref{thm:delta-method} with $\sigma^2(\fprate) =
    \fprate(1 - \fprate) / n_0$ (the variance of $\hat{\fprate}$
    from \Cref{prop:binomial-variance}). Derive
    $\Var[g(\hat{\fprate})]$ in terms of $\fprate$ and $n_0$,
    simplifying as far as possible.
  \item[(c)] In one sentence, identify the geometric feature of the
    result that makes the log-odds scale attractive for
    confidence intervals near the boundaries $\fprate \to 0$ or
    $\fprate \to 1$: the variance on the log-odds scale changes
    smoothly across the unit interval, whereas the variance of
    $\hat{\fprate}$ itself goes to zero at the endpoints.
\end{enumerate}

\item \exstarstar
\label{ex:ch7-6}
\textbf{Prevalence-dependence visualization.}
Fix the symmetric classifier $\fprate = 0.01, \fnrate = 0.01$.
\begin{enumerate}
  \item[(a)] Plot or sketch $\PPV(\pi)$ from \Cref{def:ppv} as a
    function of $\pi \in [0, 1]$. Identify the value of $\pi$ at
    which $\PPV$ is minimized, the value at which it is maximized,
    and the corresponding extremal values of $\PPV$.
  \item[(b)] Plot or sketch $\NPV(\pi)$ from \Cref{def:npv} as a
    function of $\pi \in [0, 1]$. Identify the value of $\pi$ at
    which $\NPV$ is minimized and the value at which it is
    maximized.
  \item[(c)] In one sentence, describe the qualitative symmetry
    between the two curves: $\PPV$ rises from $0$ at $\pi = 0$ to
    $1$ at $\pi = 1$, and $\NPV$ falls from $1$ at $\pi = 0$ to
    $0$ at $\pi = 1$, because $\PPV$ is a column-conditional rate
    on the positive-output column and $\NPV$ is its mirror on the
    negative-output column.
\end{enumerate}

\item \exstarstar
\label{ex:ch7-7}
\textbf{Sample-size formula for PPV CI.}
Suppose the workload prevalence $\pi$ is known and the goal is a
Wald-style 95\% confidence interval for $\PPV$ whose half-width is
at most $h$.
\begin{enumerate}
  \item[(a)] Starting from \Cref{prop:ppv-variance} and the Wald
    form $\widehat{\PPV} \pm 1.96 \sqrt{\Var[\widehat{\PPV}]}$,
    set the half-width to $h$ and solve for the total sample
    size $n$. Use $n_\top = n \pi$ and $n_\bot = n (1 - \pi)$ in
    the variance formula.
  \item[(b)] State the resulting formula for $n$ as a function of
    $\fprate, \fnrate, \pi, h$. Comment in one sentence on how
    $n$ scales with $h$ (the standard $1/h^2$ scaling shared by
    every Wald-style sample-size calculation) and on whether $n$
    grows or shrinks as $\pi \to 0$ at fixed
    $(\fprate, \fnrate, h)$.
  \item[(c)] Substitute the numbers $\fprate = 0.05, \fnrate =
    0.10, \pi = 0.1, h = 0.01$ and report the required $n$. Round
    up to the nearest integer.
\end{enumerate}

\item \exstarstar
\label{ex:ch7-8}
\textbf{Bloom filter PPV with uncertain prevalence.}
A Bloom filter has known $\fprate = 0.005$ and $\fnrate = 0$. It is
deployed on a workload whose prevalence is estimated as
$\hat\pi = 0.05$ from a Bernoulli sample of $1000$ messages. The
Wilson 95\% CI for $\pi$ is approximately $[0.038, 0.066]$.
\begin{enumerate}
  \item[(a)] Apply \Cref{def:ppv} at the point estimates
    $\fprate = 0.005, \fnrate = 0, \hat\pi = 0.05$ to compute the
    point estimate $\widehat{\PPV}$.
  \item[(b)] Treat $\PPV$ as a function of two random quantities,
    $\hat\fprate$ and $\hat\pi$, with $\hat\fprate$ holding its
    deployed value $0.005$ and $\hat\pi$ taking its empirical
    value with the stated Wilson CI. Apply the bivariate
    \Cref{thm:delta-method} (or substitute the Wilson endpoints
    directly into \Cref{def:ppv} as a quick approximation) to
    propagate the prevalence uncertainty into a CI for $\PPV$.
  \item[(c)] Report the resulting $\PPV$ interval to three decimal
    places. In one sentence, compare its width to the width that
    would be obtained by treating $\hat\pi$ as known: the second
    source of uncertainty has now entered, and the CI is
    correspondingly wider than the
    \Cref{prop:ppv-variance}-based CI computed at fixed $\pi$.
\end{enumerate}

% -----------------------------------------------------------------------
% CHALLENGING / EXTENSION  (★★★ × 3): exercises 9, 10, 11
% -----------------------------------------------------------------------

\item \exstarstarstar
\label{ex:ch7-9}
\textbf{MLE vs method-of-moments.}
\Cref{prop:binomial-mle} identified $\hat{\fprate} = \mathrm{FP}/n_0$
as the maximum-likelihood estimator. The method of moments offers an
alternative recipe: equate sample moments to population moments and
solve.
\begin{enumerate}
  \item[(a)] The first-moment method-of-moments estimator equates
    $\mathbb{E}[\mathrm{FP}] = n_0 \fprate$ to the observed
    $\mathrm{FP}$ and solves $\hat{\fprate}_{\mathrm{MoM},1} =
    \mathrm{FP}/n_0$. Confirm that this coincides with the MLE of
    \Cref{prop:binomial-mle}.
  \item[(b)] A second-moment method-of-moments estimator equates
    $\mathbb{E}[\mathrm{FP}^2] = n_0 \fprate (1 - \fprate) + n_0^2
    \fprate^2$ to $\mathrm{FP}^2$ and solves the resulting
    quadratic for $\fprate$. Carry out the algebra and confirm
    that the positive root reduces to $\mathrm{FP}/n_0$ in the
    binomial model (the two recipes coincide because the
    Bernoulli/binomial family has $\mathrm{Var}(X) = p(1-p)$
    fully determined by the mean).
  \item[(c)] Construct a small example where the two recipes
    disagree. Take a two-parameter model such as $\mathrm{Beta}(a,
    b)$ with $a, b$ both unknown: the first-moment equation alone
    is one equation in two unknowns, so the moment hierarchy must
    include the second moment to identify both parameters. State
    in one sentence the moral that the MLE = first-moment-MoM
    coincidence is special to one-parameter exponential families
    of which the binomial is the canonical example.
\end{enumerate}

\item \exstarstarstar
\label{ex:ch7-10}
\textbf{Coverage of the delta-method CI.}
The §7.4 delta-method CI for $\widehat{\PPV}$ inherits its
calibration from the underlying Wald CIs on $\hat{\fprate}$ and
$\hat{\fnrate}$ via \Cref{prop:ppv-variance}.
\begin{enumerate}
  \item[(a)] Sketch why the empirical coverage of the delta-method
    PPV CI can drop below nominal $1 - \alpha$ when the
    prevalence $\pi$ is small. Identify the term in
    \Cref{prop:ppv-variance} that becomes the dominant source of
    variability in the small-$\pi$ regime, and argue that this
    term is the false-positive count $\mathrm{FP}$ on the
    truth-$\bot$ slice, whose own Wald CI is unreliable at small
    $\fprate$ by §7.3.
  \item[(b)] Connect the failure to the §7.3 Wald sawtooth: the
    PPV CI is a ratio with a small denominator $\bar t + \bar f$
    when $\pi$ is small and the per-event $\fprate$ is also
    small, so a discrete jump in $\mathrm{FP}$ moves
    $\widehat{\PPV}$ by a substantial fraction of the CI
    half-width, replicating the sawtooth in a derived quantity.
  \item[(c)] In one or two sentences, propose a remedy: replace
    the Wald-style delta-method CI with a Wilson-style score
    interval on the column-conditional binomial
    $\TP \mid \TP + \FP \sim \mathrm{Binomial}(\TP + \FP, \PPV)$,
    which inherits the boundary correction of
    \Cref{def:wilson-ci} and stays inside $[0, 1]$ by
    construction.
\end{enumerate}

\item \exstarstarstar
\label{ex:ch7-11}
\textbf{$\Fone$ vs precision-recall tradeoff.}
The harmonic mean $\Fone = 2 \Precision \cdot \Recall / (\Precision
+ \Recall)$ of \Cref{def:f1} treats $\Precision$ and $\Recall$ on
equal footing, but the resulting iso-$\Fone$ contours in the
$(\Precision, \Recall)$ plane have geometric structure worth naming.
\begin{enumerate}
  \item[(a)] Fix the target $\Fone = 0.9$. Solve
    $0.9 = 2 \Precision \cdot \Recall / (\Precision + \Recall)$
    for $\Recall$ as a function of $\Precision$. Verify that the
    resulting expression is well-defined precisely when
    $\Precision > 0.45$ (so that the denominator stays positive
    and $\Recall \in (0, 1]$).
  \item[(b)] Plot or sketch the curve in the $(\Precision,
    \Recall)$ unit square. Identify its symmetry under
    interchange of the two axes (the curve is invariant under
    $\Precision \leftrightarrow \Recall$, a direct consequence
    of the symmetry of the harmonic mean in its two arguments)
    and the point at which $\Precision = \Recall$ on the curve.
  \item[(c)] Argue that $\Fone$ weights $\Precision$ and
    $\Recall$ equally, in the precise sense that
    $\partial \Fone / \partial \Precision = \partial \Fone /
    \partial \Recall$ at the symmetric point $\Precision =
    \Recall$. Then argue informally that any weighted
    generalization $F_\beta = (1 + \beta^2) \Precision \cdot
    \Recall / (\beta^2 \Precision + \Recall)$ with $\beta \neq
    1$ would break this symmetry, reweighting the two component
    rates by $\beta^2$ on the precision side.
\end{enumerate}

\end{enumerate}
